% Pursuing Stacks (cotes 134-2 to 134-8): archive cote 134-5, site page 23 (PDF page 24), typescript page 210.
% typescript. Draft compared: ps4.tex:1043-1101.
% AI-transcribed by Claude Opus 5.5 (claude-opus-5-5, adaptive thinking, effort high), Message Batches API, batch ps2 (msgbatch_01NVwr1Yuo8qJf5rykeGBXPK), 2026-11-12; not yet checked by eye.
\scanpage{134-5/23}
\mspage{210}

seems to be irrelevant.  It might be worth while to write down with care
what exactly is needed,  in order to define,  in terms of a bunch of
cellular structures\xout{,} of spheres,  a corresponding test-category (hopefully
even a strict one),  and a way of computing in terms of boundary operators
the cohomology of the corresponding models.  The more delicate point here
may be that if we really want to get an actual test category,  not just a
\xout{strict} weak one,  we should have \inserted[typed]{``enough''} degeneracy maps between our cellular struc-%
tures,  which might well prove an extremely stringent requirement.  Again,
it will be interesting if we can meet it otherwise than just sticking to
our trilogy.

D) ``Back to topoi''.  They have been my main intuitive leading thread in
the reflections so far,  but have remained somewhat implicit most times.
When working with categories  $X$ as models for homotopy types,  we have been
thinking in reality of the associated topos $X\hat{}$.  In the same way,  when re-%
lativizing the \deleted{(Hot)} construction \inserted{of (Hot)} as a derived category over an object $I$
of (Cat)\hand{,} \margin{the words ``the construction of (Hot) as a derived category'' on the line above are circled by hand, with an arrow pointing to this place} namely working with categories over $I$ as models for a derived
category $D(I)$,  the leading intuition again has been to look at $X$ over $I$
as the topos $X\hat{}$\hand{/}over $I\hat{}$.  This gives strong suggestions as to defining
$D(I)$ not only in terms of a given category $I$ ,  but also in terms of an
arbitrary topos $T$  as well,  standing for $I\hat{}$,  and look to what extent the
$f^{*},f_{\hand{!}},f_{*}$  formalism of derivators extends to the case when $f$ is a ``map''
between topoi.  The definition of \inserted[typed]{some} $D(T)$ for a topos $T$ can be given in a
number of ways it seems.  As far as I know,  the only one which has been
\inserted{written and} used so far consists in taking \xout{the} a suitable derived category of the cate-%
gory of semisimplicial objects of $T$ .  This is done in Illusie's thesis,
where there is no mention though of $f_{\hand{!}}$ or $f_{*}$ functors  -  which,  maybe,
should rather be written $Lf_{\hand{!}}$  and $Rf_{*}$ ,  suggesting they are respectively
left and right derived functors of the more familiar $f_{\hand{!}}$ and $f_{*}$  functors
for sheaves (\inserted{the} left and right adjoints of the inverse image functor $f^{*}$  for
sheaves of sets).  One would expect\hand{,} at  best\hand{,} $Lf_{\hand{!}}$  to exist when $f_{\hand{!}}$ itself
does,  namely \inserted{when} $f^{*}$  commutes not only with finite inverse limits,  but with
infinite products as well.  As for $Rf_{*}$ ,  whereas there is no problem for
the existence of $f_{*}$ itself,  \inserted[typed]{already} \xout{\unsure{still even}} in the case of a morphism of topoi
coming from a map in (Cat),  a map say between finite ordered sets,  the
existence of $Rf_{*}$  has still to be established.  Thus presumably I'll con-%
tent myself with writing down and comparing a few tentative definitions of
$D(T)$ and make some reasonable guesses as to it's \xout{functorial} variances.
Intuitively,  the objects of $D(T)$  may be viewed as \hand{``}sheaves of homotopy
types over $T$\hand{''} , or \hand{``}relative homotopy types over $T$\hand{''},  or \hand{``}non commutative
chain complexes over $T$ up to quasi-isomorphism\hand{''}.  As in the case when $T$ is

\margin{23}
